\chapter{Fractions et puissances}\label{ch:g9:fractions}

\omterm{def:g9:fractions:fraction}{Fractions} et \omterm{def:g9:fractions:power}{puissances} sont les outils quotidiens du calcul~: partager
des quantités, comparer des proportions, écrire des nombres très grands
ou très petits. Ce chapitre consolide les règles pour calculer avec
elles --- avec chaque étape intermédiaire écrite --- et introduit la
\omterm{def:g9:fractions:scientific}{notation scientifique}.

\section{Calculer avec des fractions}

\begin{definition}[Fraction]\label{def:g9:fractions:fraction}
Une \emph{fraction}\index{fraction} $\dfrac ab$ (avec $b \neq 0$) est le
\omterm{def:g3:division:remainder}{quotient} de $a$ par $b$. Deux fractions sont égales lorsque l'une
s'obtient de l'autre en multipliant (ou divisant) \omterm{def:g4:fractions:def}{numérateur} et
\omterm{def:g4:fractions:def}{dénominateur} par le même nombre non nul~:
\[
\frac ab = \frac{a \times k}{b \times k} \qquad (k \neq 0).
\]
\end{definition}

\begin{example}\label{ex:g9:fractions:simplify}
Simplifier $\dfrac{42}{56}$ étape par étape~:
\[
\frac{42}{56} = \frac{42 \div 2}{56 \div 2} = \frac{21}{28}
= \frac{21 \div 7}{28 \div 7} = \frac34 .
\]
(Au \cref{ch:g9:arith} le plus grand commun \omterm{def:g5:division:multiple}{diviseur} fera cela en une
seule étape.)
\end{example}

\begin{method}[Additionner et soustraire des fractions]\label{met:g9:fractions:add}
\begin{enumerate}
  \item Mettre les deux \omterm{def:g9:fractions:fraction}{fractions} sur un \emph{\omterm{def:g4:fractions:def}{dénominateur} commun}
        (un \omterm{def:g5:division:multiple}{multiple} commun des deux \omterm{def:g4:fractions:def}{dénominateurs})~;
  \item additionner ou soustraire les \omterm{def:g4:fractions:def}{numérateurs}, en gardant le
        \omterm{def:g4:fractions:def}{dénominateur}~;
  \item simplifier le résultat si possible.
\end{enumerate}
\end{method}

\begin{example}\label{ex:g9:fractions:add}
Calculer $\dfrac56 + \dfrac34$. Un \omterm{def:g4:fractions:def}{dénominateur} commun de $6$ et $4$
est $12$~:
\[
\frac56 + \frac34
= \frac{5 \times 2}{6 \times 2} + \frac{3 \times 3}{4 \times 3}
= \frac{10}{12} + \frac{9}{12}
= \frac{19}{12}.
\]
Calculer $2 - \dfrac37$. Écrire $2$ comme \omterm{def:g9:fractions:fraction}{fraction} sur $7$~:
\[
2 - \frac37 = \frac{14}{7} - \frac37 = \frac{11}{7}.
\]
\end{example}

\begin{omfigure}
\begin{tikzpicture}[scale=0.55]
  \foreach \i in {0,...,11}
    \draw[omExo] ({mod(\i,12)},0) rectangle ({mod(\i,12)+1},1);
  \foreach \i in {0,...,9}
    \fill[omDef!35] (\i,0) rectangle (\i+1,1);
  \draw[thick] (0,0) rectangle (12,1);
  \node[below, font=\small] at (6,-0.1)
    {$\frac56 = \frac{10}{12}$~: 10 cases sur 12};
  \begin{scope}[yshift=-2.4cm]
  \foreach \i in {0,...,11}
    \draw[omExo] (\i,0) rectangle (\i+1,1);
  \foreach \i in {0,...,8}
    \fill[omThm!35] (\i,0) rectangle (\i+1,1);
  \draw[thick] (0,0) rectangle (12,1);
  \node[below, font=\small] at (6,-0.1)
    {$\frac34 = \frac{9}{12}$~: 9 cases sur 12};
  \end{scope}
\end{tikzpicture}

{\small Pourquoi un \omterm{def:g4:fractions:def}{dénominateur} commun~: découper les deux entiers en
$12$ cases égales rend les \omterm{def:g9:fractions:fraction}{fractions} comparables et additionnables ---
ensemble, $10 + 9 = 19$ douzièmes.}
\end{omfigure}

\begin{method}[Multiplier et diviser des fractions]\label{met:g9:fractions:mult}
\begin{enumerate}
  \item Pour multiplier, multiplier les \omterm{def:g4:fractions:def}{numérateurs} entre eux et les
        \omterm{def:g4:fractions:def}{dénominateurs} entre eux~: $\dfrac ab \times \dfrac cd =
        \dfrac{ac}{bd}$ (simplifier \emph{avant} de multiplier dès que
        possible)~;
  \item pour diviser, multiplier par l'\omterm{def:g8:fractions:inverse}{inverse} du \omterm{def:g5:division:multiple}{diviseur}~:
        $\dfrac ab \div \dfrac cd = \dfrac ab \times \dfrac dc$.
\end{enumerate}
\end{method}

\begin{example}\label{ex:g9:fractions:mult}
\[
\frac{7}{15} \times \frac{25}{14}
= \frac{7 \times 25}{15 \times 14}
= \frac{7 \times 25}{15 \times 2 \times 7}
= \frac{25}{30} = \frac56,
\]
où l'on a simplifié par le facteur commun $7$, puis par $5$. Et une
division~:
\[
\frac{3}{4} \div \frac{9}{8}
= \frac34 \times \frac89
= \frac{3 \times 8}{4 \times 9}
= \frac{24}{36} = \frac23 .
\]
\end{example}

\section{Puissances}

\begin{definition}[Puissance]\label{def:g9:fractions:power}
Pour un nombre réel $a$ et un entier positif $n$, la
\emph{puissance}\index{puissance} $a^n$ est le \omterm{def:g2:mult:def}{produit} de $n$ facteurs
égaux à $a$~:
\[
a^n = \underbrace{a \times a \times \dots \times a}_{n \text{ facteurs}} .
\]
Par convention $a^0 = 1$ (pour $a \neq 0$), et les \omterm{def:g8:powers:def}{exposants} négatifs
désignent les \omterm{def:g8:fractions:inverse}{inverses}~:
\[
a^{-n} = \frac{1}{a^n}.
\]
\end{definition}

\begin{example}
$2^4 = 2 \times 2 \times 2 \times 2 = 16$~; \quad
$(-3)^2 = 9$ mais $-3^2 = -(3^2) = -9$ (l'\omterm{def:g8:powers:def}{exposant} lie avant le signe
moins)~; \quad $10^{-3} = \frac{1}{1000} = 0.001$~; \quad
$5^1 = 5$ et $5^0 = 1$.
\end{example}

\begin{theorem}[Règles des exposants]\label{thm:g9:fractions:rules}
Pour tous $a$, $b$ non nuls et tous entiers $m$, $n$~:
\[
a^m \times a^n = a^{m+n},
\qquad
\frac{a^m}{a^n} = a^{m-n},
\qquad
\left(a^m\right)^n = a^{mn},
\qquad
(ab)^n = a^n b^n .
\]
\end{theorem}

\begin{proof}
Pour les \omterm{def:g8:powers:def}{exposants} positifs, compter les facteurs. Pour la première
règle~: $a^m \times a^n$ est $m$ facteurs $a$ suivis de $n$ facteurs
$a$, au total $m + n$ facteurs. Pour la troisième~:
$\left(a^m\right)^n$ répète $n$ fois un bloc de $m$ facteurs, donnant
$mn$ facteurs. Les deux autres règles sont analogues~; l'extension aux
\omterm{def:g8:powers:def}{exposants} nuls et négatifs se vérifie à partir de
$a^{-n} = \frac{1}{a^n}$ (et rend la convention $a^0 = 1$ le seul choix
cohérent, puisque $a^n \times a^0$ doit valoir $a^{n+0} = a^n$).
\end{proof}

\begin{example}\label{ex:g9:fractions:rules}
Simplifier étape par étape~:
\[
\frac{3^7 \times 3^{-2}}{3^4}
= \frac{3^{7 + (-2)}}{3^4}
= \frac{3^5}{3^4}
= 3^{5-4} = 3^1 = 3 .
\]
Et avec deux lettres~: $(2a^3)^2 \times a = 4 a^6 \times a = 4a^7$.
\end{example}

\section{Puissances de dix et notation scientifique}

\begin{proposition}[Puissances de dix]\label{prop:g9:fractions:ten}
Pour un entier positif $n$~: $10^n$ s'écrit «~$1$ suivi de $n$ \omterm{def:g1:counting:zero}{zéros}~»,
et $10^{-n} = 0.0\dots01$ avec le $1$ à la $n$-ième place décimale.
Multiplier un \omterm{def:g6:decimals:places}{nombre décimal} par $10^n$ (resp.\ $10^{-n}$) décale sa
\omterm{def:g4:decimals:point}{virgule} de $n$ places vers la droite (resp.\ la gauche).
\end{proposition}

\begin{definition}[Notation scientifique]\label{def:g9:fractions:scientific}
La \emph{notation scientifique}\index{notation scientifique} d'un
\omterm{def:g6:decimals:places}{nombre décimal} positif est la façon unique de l'écrire sous la forme
\[
a \times 10^n,
\qquad\text{avec } 1 \leq a < 10 \text{ et } n \text{ un entier.}
\]
\end{definition}

\begin{example}\label{ex:g9:fractions:scientific}
La distance Terre--Soleil est d'environ $149\,600\,000$ km $= 1.496
\times 10^8$ km. La taille d'une molécule d'eau est d'environ
$0.000\,000\,000\,28$ m $= 2.8 \times 10^{-10}$ m. Pour calculer avec de
tels nombres, regrouper les parties décimales et les \omterm{def:g9:fractions:power}{puissances} de
dix~:
\[
(3 \times 10^5) \times (2.5 \times 10^{-2})
= (3 \times 2.5) \times 10^{5 + (-2)}
= 7.5 \times 10^{3}.
\]
\end{example}

\begin{method}[Mettre un nombre en notation scientifique]\label{met:g9:fractions:scientific}
\begin{enumerate}
  \item Déplacer la \omterm{def:g4:decimals:point}{virgule} juste après le premier chiffre non nul~;
        cela donne le facteur $a$ avec $1 \leq a < 10$~;
  \item compter le nombre de places dont la \omterm{def:g4:decimals:point}{virgule} a bougé~: ce
        décompte est $n$, positif si le nombre d'origine était
        $\geq 10$, négatif s'il était $< 1$~;
  \item vérifier~: $a \times 10^n$ doit reproduire le nombre d'origine.
\end{enumerate}
\end{method}

\section{Exercices}

\begin{exercise}[$\star$]\label{exo:g9:fractions:1}
Calculer et donner le résultat sous forme de \omterm{def:g9:fractions:fraction}{fraction} entièrement
simplifiée~:
\[
\frac13 + \frac15, \qquad
\frac72 - \frac53, \qquad
\frac{5}{12} + \frac38, \qquad
3 - \frac45 .
\]
\end{exercise}

\begin{exercise}[$\star$]\label{exo:g9:fractions:2}
Calculer et simplifier~:
\[
\frac59 \times \frac{27}{35}, \qquad
\frac{8}{15} \div \frac{4}{5}, \qquad
\frac23 \times \frac34 \times \frac45 .
\]
\end{exercise}

\begin{exercise}[$\star$]\label{exo:g9:fractions:3}
Calculer~:
\[
2^5, \qquad (-2)^4, \qquad -2^4, \qquad 4^{-2}, \qquad
\left(\tfrac23\right)^3, \qquad 7^0 .
\]
\end{exercise}

\begin{exercise}[$\star$]\label{exo:g9:fractions:4}
Écrire comme une seule \omterm{def:g9:fractions:power}{puissance}~:
\[
5^3 \times 5^6, \qquad
\frac{7^9}{7^5}, \qquad
\left(2^3\right)^4, \qquad
\frac{3^2 \times 3^7}{3^5}, \qquad
2^5 \times 4 .
\]
\end{exercise}

\begin{exercise}[$\star$]\label{exo:g9:fractions:5}
Écrire en \omterm{def:g9:fractions:scientific}{notation scientifique}~: $45\,000$~; $0.0072$~; $302.5$~;
$0.000\,000\,9$~; douze millions.
\end{exercise}

\begin{exercise}[$\star\star$]\label{exo:g9:fractions:6}
Calculer, en donnant le résultat en \omterm{def:g9:fractions:scientific}{notation scientifique}~:
\[
(2 \times 10^7) \times (4.5 \times 10^{-3}),
\qquad
\frac{9 \times 10^5}{3 \times 10^{-2}},
\qquad
(5 \times 10^3)^2 .
\]
\end{exercise}

\begin{exercise}[$\star\star$]\label{exo:g9:fractions:7}
Un réservoir est plein aux $\frac35$. Après avoir ajouté $30$ litres,
il est plein aux $\frac34$. Quelle est la capacité du réservoir~?
(Exprimer d'abord la \omterm{def:g9:fractions:fraction}{fraction} ajoutée du réservoir.)
\end{exercise}

\begin{exercise}[$\star\star$]\label{exo:g9:fractions:8}
Alice mange $\frac14$ d'un gâteau, puis Ben mange $\frac13$ de ce qui
reste, puis Carla mange $\frac12$ de ce qui reste encore. Quelle
\omterm{def:g9:fractions:fraction}{fraction} du gâteau reste-t-il à la fin~? (Calculer le \omterm{def:g3:division:remainder}{reste} après
chaque étape.)
\end{exercise}

\begin{exercise}[$\star\star$]\label{exo:g9:fractions:9}
La lumière parcourt $3 \times 10^8$ mètres par seconde. Le Soleil est
à environ $1.5 \times 10^{11}$ m. Combien de temps la lumière du Soleil
met-elle à nous atteindre~? Donner la réponse en secondes (\omterm{def:g9:fractions:scientific}{notation
scientifique} non nécessaire), puis en minutes et secondes.
\end{exercise}

\begin{exercise}[$\star\star\star$]\label{exo:g9:fractions:10}
Lequel est plus grand, $2^{30}$ ou $3^{20}$~? (Indication~: écrire les
deux comme des \omterm{def:g9:fractions:power}{puissances} d'\omterm{def:g8:powers:def}{exposant} $10$ en utilisant
$\left(a^m\right)^n = a^{mn}$, et comparer $2^3$ avec $3^2$.)
\end{exercise}

\section{Problème~: le code secret des décimales périodiques}

\begin{problem}[{Devoir maison --- toute fraction a une écriture
décimale qui s'arrête ou se répète, tout décimal périodique est une
fraction, et $0.999\ldots$ vaut exactement $1$}]
\label{pb:g9:fractions:1}
Divisons $3$ par $11$~: les chiffres se mettent à psalmodier,
$0.272727\ldots$ sans fin. Divisons $1$ par $4$~: les chiffres
s'arrêtent net, $0.25$. Ce problème démontre que ce sont là les deux
seuls comportements possibles pour une \omterm{def:g9:fractions:fraction}{fraction} --- puis il \omterm{def:g8:fractions:inverse}{inverse}
la machine~: étant donné un décimal périodique quelconque, il
reconstitue la \omterm{def:g9:fractions:fraction}{fraction} qui se cache derrière. En chemin, il règle
une fois pour toutes l'égalité la plus discutée des mathématiques
scolaires.

\medskip
\noindent\textbf{Partie I --- De la \omterm{def:g9:fractions:fraction}{fraction} au décimal.}
\begin{enumerate}
  \item Calculer, en posant la division, les écritures décimales de
        $\frac14$, $\frac13$, $\frac56$ et $\frac{3}{11}$.
        Lesquelles s'arrêtent, lesquelles se répètent, et avec
        quelle période~?
  \item Expliquer pourquoi l'écriture décimale de \emph{n'importe
        quelle} \omterm{def:g9:fractions:fraction}{fraction} $\frac ab$ doit soit s'arrêter, soit finir
        par se répéter. (Dans la division posée par $b$, quelles
        valeurs les \omterm{def:g3:division:remainder}{restes} successifs peuvent-ils prendre~? Que se
        passe-t-il à l'instant où un \omterm{def:g3:division:remainder}{reste} apparaît pour la seconde
        fois~?)
  \item Pousser la division de $1$ par $7$ assez loin pour trouver
        la période. Quelle est sa longueur, et comment cette
        longueur se compare-t-elle à la borne promise par la
        question 2~?
  \item Certaines \omterm{def:g9:fractions:fraction}{fractions} s'arrêtent parce que leur \omterm{def:g4:fractions:def}{dénominateur}
        peut être gonflé jusqu'à une \omterm{def:g9:fractions:power}{puissance} de dix~:
        $\frac14 = \frac{25}{100}$,
        $\frac{3}{8} = \frac{375}{1000}$. Expliquer le critère, et
        pourquoi aucun multiplicateur entier ne pourra jamais
        transformer $3$, $6$ ou $11$ en $10$, $100$,
        $1000, \dots$ (par quel chiffre le \omterm{def:g2:mult:def}{produit}
        finirait-il~?).
  \item Sans poser de division, ranger $\frac{9}{40}$,
        $\frac{5}{12}$ et $\frac{33}{50}$ en «~s'arrête~» et
        «~se répète~», puis donner l'écriture décimale des deux qui
        s'arrêtent.
\end{enumerate}

\medskip
\noindent\textbf{Partie II --- Du décimal périodique à la \omterm{def:g9:fractions:fraction}{fraction}.}
\begin{enumerate}[resume]
  \item L'astuce du décalage. Posons $x = 0.7777\ldots$ Calculer
        $10x$, puis $10x - x$, et en déduire $x$ sous forme de
        \omterm{def:g9:fractions:fraction}{fraction}. Vérifier en posant la division.
  \item Posons $x = 0.727272\ldots$ Quelle \omterm{def:g9:fractions:power}{puissance} de dix décale
        $x$ d'exactement une période~? S'en servir pour écrire $x$
        sous forme de \omterm{def:g9:fractions:fraction}{fraction} irréductible.
  \item Un cas mixte~: $x = 0.58333\ldots$ (seul le $3$ se répète).
        Calculer $1000x - 100x$ et en déduire $x$ sous forme de
        \omterm{def:g9:fractions:fraction}{fraction} irréductible.
  \item La célèbre~: appliquer l'astuce du décalage à
        $x = 0.9999\ldots$ Quelle \omterm{def:g9:fractions:fraction}{fraction} --- quel \emph{nombre}
        --- trouve-t-on~? Réconcilier le résultat avec son
        intuition, en se souvenant du voisin fantôme du
        \cref{pb:g6:decimals:1} et du \omterm{def:g3:division:remainder}{reste} qui rétrécit du
        \cref{pb:g6:fractions:1}~: $0.999\ldots$ est-il «~juste en
        dessous~» de $1$, ou bien un autre nom pour $1$~?
  \item Convertir $2.454545\ldots$ et $0.123123123\ldots$ en
        \omterm{def:g9:fractions:fraction}{fractions} irréductibles.
\end{enumerate}

\medskip
\noindent\textbf{Partie III --- Le dictionnaire achevé.}
\begin{enumerate}[resume]
  \item Les questions 6 à 10 suggèrent un dictionnaire~: une
        période de $k$ chiffres démarrant juste après la \omterm{def:g4:decimals:point}{virgule}
        vaut cette période sur $k$ neufs
        ($0.727272\ldots = \frac{72}{99}$). Utiliser le
        dictionnaire pour prévoir $0.037037\ldots$ sous forme de
        \omterm{def:g9:fractions:fraction}{fraction}, simplifier --- et s'émerveiller~: quelle \omterm{def:g9:fractions:fraction}{fraction}
        étonnamment simple psalmodie «~$037$~»~? Le vérifier en
        posant la division.
  \item Combiner le dictionnaire avec un décalage~: écrire
        $0.0272727\ldots$ sous forme de \omterm{def:g9:fractions:fraction}{fraction} irréductible.
  \item Le nombre $0.101001000100001\ldots$ (un $1$, puis un $0$,
        puis un $1$, puis deux $0$, puis un $1$, puis trois $0$, et
        ainsi de suite) ne s'arrête jamais. Est-ce un décimal
        périodique~? Que dit alors la partie I à son sujet ---
        peut-il être une \omterm{def:g9:fractions:fraction}{fraction}~? (On vient de rencontrer son
        premier nombre démontrablement \emph{irrationnel}~; le plus
        célèbre est démasqué au \cref{ch:g9:sqrt}.)
  \item Combien de chiffres a $2^{30}$, en gros~? Utiliser
        $2^{10} = 1\,024 \approx 1.024 \times 10^3$ et les règles
        des \omterm{def:g8:powers:def}{exposants} (\cref{thm:g9:fractions:rules}) pour écrire
        $2^{30}$ en \omterm{def:g9:fractions:scientific}{notation scientifique} (approchée), et conclure.
        (Comparer avec \cref{exo:g9:fractions:10}.)
  \item Bouquet final~: écrire $0.20262026\ldots$ sous forme de
        \omterm{def:g9:fractions:fraction}{fraction}, et énoncer la morale à double sens de tout le
        problème~: quelles écritures décimales sont des \omterm{def:g9:fractions:fraction}{fractions},
        et quelles \omterm{def:g9:fractions:fraction}{fractions} ont quelles écritures décimales~?
\end{enumerate}
\end{problem}
